% =============================================================================
% template.tex — КиТГ (Комбинаторика и Теория Графов)
% Лектор: Казакевич В.Г. | Конспект: Котельников
% =============================================================================

% --- Кодировка и язык ---
\usepackage[utf8]{inputenc}
\DeclareUnicodeCharacter{03B1}{\ensuremath{\alpha}}
\DeclareUnicodeCharacter{03B2}{\ensuremath{\beta}}
\DeclareUnicodeCharacter{03B3}{\ensuremath{\gamma}}
\DeclareUnicodeCharacter{03B4}{\ensuremath{\delta}}
\DeclareUnicodeCharacter{03B5}{\ensuremath{\varepsilon}}
\DeclareUnicodeCharacter{03B6}{\ensuremath{\zeta}}
\DeclareUnicodeCharacter{03B7}{\ensuremath{\eta}}
\DeclareUnicodeCharacter{03B8}{\ensuremath{\theta}}
\DeclareUnicodeCharacter{03C3}{\ensuremath{\sigma}}
\DeclareUnicodeCharacter{03C7}{\ensuremath{\chi}}
\usepackage[T2A]{fontenc}
\usepackage[english,russian]{babel}

% --- Геометрия ---
\usepackage[margin=1.5cm]{geometry}

% --- Математика ---
\usepackage{amsmath}
\usepackage{amssymb}
\usepackage{amsthm}
\usepackage{mathtools}

% --- Типографика ---
\usepackage{microtype}
\usepackage{indentfirst}

% --- Списки ---
\usepackage{enumitem}

% --- Графика ---
\usepackage{graphicx}
\usepackage{needspace}

% --- Единый размер картинок-графов (нормализация) ---
% \grfig{path} — центрированная картинка стандартной ширины (0.5\textwidth).
% \grfig[0.7]{path} — задать другую ширину при необходимости.
\newcounter{grfignum}
\newcommand{\grfig}[3][0.5]{%
  \begin{center}%
  \includegraphics[width=#1\textwidth]{#2}\\[3pt]%
  \refstepcounter{grfignum}{\small\textit{Рис.~\thegrfignum. #3}}%
  \end{center}}
\newcommand{\grcap}[1]{\begin{center}\refstepcounter{grfignum}%
  {\small\textit{Рис.~\thegrfignum. #1}}\end{center}}

% --- \exhead{Пример 1.} — заголовок примера, который не отрывается от
%     своего содержимого (рисунка/таблицы) при переносе страниц.
\newcommand{\exhead}[1]{\par\needspace{12\baselineskip}%
  \medskip\noindent\textbf{#1}\nopagebreak\hspace{0.33em}\ignorespaces}
\usepackage{caption}
\usepackage{subcaption}

% --- Таблицы ---
\usepackage{booktabs}

% --- TikZ / PGFPlots ---
\usepackage{tikz}
\usepackage{pgfplots}
\pgfplotsset{compat=1.18}
\usetikzlibrary{arrows.meta, positioning, calc}

% --- Цвета ---
\usepackage{xcolor}

% Цвета для окружений
\definecolor{defngreen}{HTML}{388E3C}
\definecolor{stmtblue}{HTML}{1565C0}
\definecolor{lemmablue}{HTML}{1976D2}
\definecolor{algopurple}{HTML}{7B1FA2}
\definecolor{taskorange}{HTML}{F57C00}
\definecolor{proofgray}{HTML}{757575}

% --- tcolorbox ---
\usepackage[most]{tcolorbox}
\tcbuselibrary{breakable, skins}

% --- Псевдокод ---
\usepackage[ruled,vlined,linesnumbered,resetcount]{algorithm2e}
\SetAlgorithmName{Алгоритм}{Алгоритм}{Список алгоритмов}
\SetKwInput{KwData}{Вход}
\SetKwInput{KwResult}{Выход}
\SetKwIF{If}{ElseIf}{Else}{если}{то}{иначе если}{иначе}{конец условия}
\SetKwFor{While}{пока}{выполнять}{конец цикла}
\SetKwFor{For}{для}{выполнять}{конец цикла}
\SetKwFor{ForEach}{для каждого}{выполнять}{конец цикла}
\SetKw{Return}{вернуть}

% --- Гиперссылки (ЗАГРУЖАТЬ ПОСЛЕДНИМ из обычных пакетов) ---
\usepackage[
  colorlinks=true,
  linkcolor=stmtblue,
  citecolor=defngreen,
  urlcolor=algopurple,
  bookmarks=true,
  bookmarksnumbered=true,
]{hyperref}

% --- cleveref (ПОСЛЕ hyperref!) ---
\usepackage[russian]{cleveref}

% --- Заголовки ---
\usepackage{titlesec}

% Форматирование глав и секций — без подписи «Глава N» (показываем только заголовок)
\titleformat{\chapter}[hang]
  {\normalfont\huge\bfseries}
  {}{0pt}{\Huge}
\titlespacing*{\chapter}{0pt}{-20pt}{20pt}


% =============================================================================
% СЧЁТЧИКИ — СКВОЗНЫЕ, НЕЗАВИСИМЫЕ (не привязаны к chapter/section)
% =============================================================================
\newcounter{defncounter}      % Определения: 1, 2, 3, ...
\newcounter{thmcounter}       % Теоремы/леммы: 1, 2, 3, ...
\newcounter{algocounter}      % Алгоритмы: 1, 2, 3, ...
\newcounter{taskcounter}      % Задачи: 1, 2, 3, ...


% =============================================================================
% ОКРУЖЕНИЯ (tcolorbox)
% Общее правило: текст внутри всех боксов — ВСЕГДА чёрный (coltext=black)
% =============================================================================

% --- 1. defn: Определение (зелёная рамка) ---
% Использование: \begin{defn}{Название}{ключ} ... \end{defn}
\newtcolorbox{defn}[2]{%
  enhanced,
  breakable,
  colframe=defngreen,
  colback=white,
  coltitle=black,
  coltext=black,
  fonttitle=\bfseries,
  title={\refstepcounter{defncounter}Определение~\thedefncounter.\ #1},
  label={def:#2},
  boxrule=1.2pt,
  arc=2pt,
  left=6pt, right=6pt, top=4pt, bottom=4pt,
  before upper={\parindent=1.5em},
}

% --- 2. statement: Теорема/Утверждение (синяя рамка) ---
% Использование: \begin{statement}{Название}{ключ} ... \end{statement}
\newtcolorbox{statement}[2]{%
  enhanced,
  breakable,
  colframe=stmtblue,
  colback=white,
  coltitle=white,
  coltext=black,
  fonttitle=\bfseries,
  title={\refstepcounter{thmcounter}Теорема~\thethmcounter.\ #1},
  label={thm:#2},
  boxrule=1.2pt,
  arc=2pt,
  left=6pt, right=6pt, top=4pt, bottom=4pt,
  before upper={\parindent=1.5em},
}

% --- 3. lemma: Лемма (синяя рамка, чуть светлее) ---
% Использует thmcounter (общий с теоремами!)
% Использование: \begin{lemma}{Название}{ключ} ... \end{lemma}
\newtcolorbox{lemma}[2]{%
  enhanced,
  breakable,
  colframe=lemmablue,
  colback=white,
  coltitle=white,
  coltext=black,
  fonttitle=\bfseries,
  title={\refstepcounter{thmcounter}Лемма~\thethmcounter.\ #1},
  label={thm:#2},
  boxrule=1.2pt,
  arc=2pt,
  left=6pt, right=6pt, top=4pt, bottom=4pt,
  before upper={\parindent=1.5em},
}

% --- 4. algo: Алгоритм (фиолетовая рамка) ---
% Использование: \begin{algo}{Название}{ключ} ... \end{algo}
\newtcolorbox{algo}[2]{%
  enhanced,
  breakable,
  colframe=algopurple,
  colback=white,
  coltitle=white,
  coltext=black,
  fonttitle=\bfseries,
  title={\refstepcounter{algocounter}Алгоритм~\thealgocounter.\ #1},
  label={algo:#2},
  boxrule=1.2pt,
  arc=2pt,
  left=6pt, right=6pt, top=4pt, bottom=4pt,
  before upper={\parindent=1.5em},
}

% --- 5. task: Задача (оранжевая рамка) ---
% Использование: \begin{task}{Название}{ключ} ... \end{task}
\newtcolorbox{task}[2]{%
  enhanced,
  breakable,
  colframe=taskorange,
  colback=white,
  coltitle=black,
  coltext=black,
  fonttitle=\bfseries,
  title={\refstepcounter{taskcounter}Задача~\thetaskcounter.\ #1},
  label={task:#2},
  boxrule=1.2pt,
  arc=2pt,
  left=6pt, right=6pt, top=4pt, bottom=4pt,
  before upper={\parindent=1.5em},
}

% --- 6. proof: Доказательство (обычный чёрный текст, жирный зачин, □ в конце) ---
\renewenvironment{proof}[1][Доказательство]{%
  \par\vspace{4pt}%
  \noindent\textbf{#1.}\ \normalfont\color{black}%
}{%
  \hfill$\square$\par\vspace{6pt}%
}


% =============================================================================
% ОКРУЖЕНИЯ-СБОРНИКИ (для summary_*.tex)
% =============================================================================

% Робастная ссылка «(см. стр. N)» — пустая, если метка не определена
% (нужно для отдельных PDF-сборников, где метки оригинала отсутствуют).
\makeatletter
\newcommand{\sumpageref}[1]{\@ifundefined{r@#1}{}{ {\normalfont\normalsize(см.~стр.~\hyperref[#1]{\pageref*{#1}})}}}
\makeatother

% summarydefn — Определение в сборнике (без инкремента счётчика)
\newtcolorbox{summarydefn}[3]{%
  enhanced,
  breakable,
  colframe=defngreen,
  colback=white,
  coltitle=black,
  coltext=black,
  fonttitle=\bfseries,
  title={Определение~#1.\ #2\sumpageref{#3}},
  boxrule=1pt,
  arc=2pt,
  left=6pt, right=6pt, top=4pt, bottom=4pt,
  before upper={\parindent=1.5em},
}

% summarythm — Теорема/Утверждение в сборнике
\newtcolorbox{summarythm}[3]{%
  enhanced,
  breakable,
  colframe=stmtblue,
  colback=white,
  coltitle=white,
  coltext=black,
  fonttitle=\bfseries,
  title={Теорема~#1.\ #2\sumpageref{#3}},
  boxrule=1pt,
  arc=2pt,
  left=6pt, right=6pt, top=4pt, bottom=4pt,
  before upper={\parindent=1.5em},
}

% summarylemma — Лемма в сборнике
\newtcolorbox{summarylemma}[3]{%
  enhanced,
  breakable,
  colframe=lemmablue,
  colback=white,
  coltitle=white,
  coltext=black,
  fonttitle=\bfseries,
  title={Лемма~#1.\ #2\sumpageref{#3}},
  boxrule=1pt,
  arc=2pt,
  left=6pt, right=6pt, top=4pt, bottom=4pt,
  before upper={\parindent=1.5em},
}

% summaryalgo — Алгоритм в сборнике
\newtcolorbox{summaryalgo}[3]{%
  enhanced,
  breakable,
  colframe=algopurple,
  colback=white,
  coltitle=white,
  coltext=black,
  fonttitle=\bfseries,
  title={Алгоритм~#1.\ #2\sumpageref{#3}},
  boxrule=1pt,
  arc=2pt,
  left=6pt, right=6pt, top=4pt, bottom=4pt,
  before upper={\parindent=1.5em},
}

% summarytask — Задача в сборнике
\newtcolorbox{summarytask}[3]{%
  enhanced,
  breakable,
  colframe=taskorange,
  colback=white,
  coltitle=black,
  coltext=black,
  fonttitle=\bfseries,
  title={Задача~#1.\ #2\sumpageref{#3}},
  boxrule=1pt,
  arc=2pt,
  left=6pt, right=6pt, top=4pt, bottom=4pt,
  before upper={\parindent=1.5em},
}


% =============================================================================
% КОМАНДЫ ДЛЯ ВСТАВКИ СБОРНИКОВ
% =============================================================================

\newcommand{\printdefinitionsummary}{%
  \IfFileExists{summary_defs.tex}{\input{summary_defs.tex}}%
    {\textit{Сборник определений пока пуст. Запустите make\_summary.py.}}%
}

\newcommand{\printtheoremsummary}{%
  \IfFileExists{summary_thms.tex}{\input{summary_thms.tex}}%
    {\textit{Сборник теорем пока пуст. Запустите make\_summary.py.}}%
}

\newcommand{\printalgorithmsummary}{%
  \IfFileExists{summary_algos.tex}{\input{summary_algos.tex}}%
    {\textit{Сборник алгоритмов пока пуст. Запустите make\_summary.py.}}%
}

\newcommand{\printtasksummary}{%
  \IfFileExists{summary_tasks.tex}{\input{summary_tasks.tex}}%
    {\textit{Сборник задач пока пуст. Запустите make\_summary.py.}}%
}


% =============================================================================
% УДОБНЫЕ МАКРОСЫ ДЛЯ КиТГ
% =============================================================================

% Множества и графы
\newcommand{\V}{\mathcal{V}}         % множество вершин
\newcommand{\E}{\mathcal{E}}         % множество рёбер
\providecommand{\G}{}
\renewcommand{\G}{\mathcal{G}}       % граф
\newcommand{\A}{\mathcal{A}}         % множество дуг (ориентированный граф)

% Операторы
\DeclareMathOperator{\dist}{dist}
\DeclareMathOperator{\diam}{diam}
\DeclareMathOperator{\rad}{rad}
\let\deg\relax
\DeclareMathOperator{\deg}{deg}
\DeclareMathOperator{\indeg}{indeg}
\DeclareMathOperator{\outdeg}{outdeg}
\DeclareMathOperator{\adj}{adj}
\DeclareMathOperator{\comp}{comp}

% Комбинаторика
\newcommand{\cmark}{\ding{51}}       % для чеклистов (нужен pifont)
\newcommand{\Nset}{\mathbb{N}}
\newcommand{\Zset}{\mathbb{Z}}
\newcommand{\Rset}{\mathbb{R}}
\newcommand{\Qset}{\mathbb{Q}}

% Кратко: множество вершин/рёбер графа G
\newcommand{\VG}{V(\G)}
\newcommand{\EG}{E(\G)}

% Делимость (рус. нотация): \vdots — «кратно», \ndiv — «не кратно»
\providecommand{\ndiv}{\mathrel{\not\vdots}}
